\section{Replication Architecture and Implementation}

\subsection{Software architecture}
\begin{draftpoints}
  \item implementation separates expandable model layers, base pretraining, incremental learning, neurogenesis control, replay generation, threshold estimation, experiment configuration, and artifact export into independently testable components
  \item \artifact{src/models/ng_linear.py} resizes dense layers while preserving mature parameters, and \artifact{src/models/ng_autoencoder.py} coordinates symmetric encoder--decoder expansion
  \item \artifact{src/training/neurogenesis_trainer.py} implements outlier selection, layer-local growth, plasticity, stability, next-layer propagation, and growth diagnostics
  \item \artifact{src/utils/intrinsic_replay.py} fits and samples class-conditional latent Gaussians, while \artifact{src/utils/dataset_replay.py} provides an original-data upper-bound control
  \item Paper configurations are stored under \artifact{config/paper/}; executable experiment, ablation, plotting, resource, and export workflows are stored under \artifact{scripts/}
\end{draftpoints}

\subsection{Replication execution pipeline}
\begin{figure}[htbp]
\centering
\resizebox{0.88\textwidth}{!}{%
\begin{tikzpicture}[
  node distance=7mm,
  box/.style={draw,rounded corners,minimum width=49mm,minimum height=8mm,align=center,font=\small},
  decision/.style={draw,diamond,aspect=2.4,align=center,font=\scriptsize,inner sep=1.5pt},
  arrow/.style={-{Latex[length=2mm]},thick},
  note/.style={font=\scriptsize,align=left,text width=38mm}
]
\node[box,fill=ProtocolBlue] (base) {Pretrain stacked autoencoder\\on base classes 1 and 7};
\node[box,below=of base] (threshold) {Estimate per-depth reconstruction thresholds\\and store replay statistics};
\node[box,below=of threshold,fill=ProtocolLightPurple] (class) {Introduce the next curriculum class};
\node[box,below=of class] (evaluate) {Compute global reconstruction error\\at every depth};
\node[decision,below=8mm of evaluate] (grow) {Outliers exceed\\allowed quota?};
\node[box,below left=8mm and 15mm of grow,fill=ProtocolRed] (local) {Add neurons; local plasticity\\and replay stabilization};
\node[box,below right=8mm and 15mm of grow,fill=ProtocolBlue] (finish) {Finalize class; refresh replay\\statistics and evaluate all classes};
\draw[arrow] (base) -- (threshold);
\draw[arrow] (threshold) -- (class);
\draw[arrow] (class) -- (evaluate);
\draw[arrow] (evaluate) -- (grow);
\draw[arrow] (grow) -- node[above left,font=\scriptsize]{yes} (local);
\draw[arrow] (grow) -- node[above right,font=\scriptsize]{no / cap} (finish);
\draw[arrow] (local.west) -| ++(-10mm,0) |- (evaluate.west);
\draw[arrow] (finish.east) -| ++(10mm,0) |- node[pos=.25,right,font=\scriptsize]{next class} (class.east);
\end{tikzpicture}
}
\caption{Implemented replication pipeline. The local growth loop repeats within each autoencoder depth until the outlier quota is satisfied or the applicable growth allowance is exhausted.}
\label{fig:replication-pipeline}
\end{figure}

\subsection{Paper-to-code fidelity}
\begin{draftpoints}
  \item implementation matches the high-level operations of Algorithm 1 after corrections to outlier identity, tensor cloning, replay isolation, threshold provenance, and experiment-regime selection
  \item Paper-oriented configurations use partial global reconstruction error for selecting incoming-class outliers, but train new features through an isolated single-hidden-layer autoencoder objective
  \item Mature encoder weights remain frozen during plasticity, decoder weights receive the reduced learning rate, and replay stabilization operates at the current level
  \item Literal equivalence cannot be claimed because several consequential choices are undocumented in the paper and must be resolved by transparent configuration
\end{draftpoints}

\begin{table}[htbp]
\centering
\caption{Condensed operation-level fidelity map; the complete audit is preserved in \artifact{docs/algorithm_fidelity.md}.}
\label{tab:fidelity-map}
\small
\begin{tabularx}{\textwidth}{p{3.9cm}p{4.4cm}p{1.9cm}Y}
\toprule
Published operation & Local implementation & Status & Consequence \\
\midrule
Global reconstruction error at each depth & Partial encode--decode forward pass and per-sample MSE & Match & Growth responds to pixel reconstruction demand \\
Select incoming-class outliers & Indexed incoming-class subset above threshold & Match & Replay samples cannot become growth candidates \\
Add and train new local nodes & Expandable linear layers plus local SHL-AE & Match & Local growth objective follows Algorithm 1 \\
Freeze mature encoder during plasticity & Gradient masks / selected parameter scope & Match & Existing feature detectors are protected \\
Decoder and stability rate $LR/100$ & Explicit learning-rate ratios & Match & Published stability asymmetry is retained \\
Class-conditional Gaussian replay & Mean, covariance factor, latent sampling, decoding & Match & No old images are reopened in IR conditions \\
Optimizer, activation, exact phase lengths & Configured, because publication is silent & Unspecified & Paper-compatible values require sensitivity analysis \\
Number of nodes added per round & Absolute or proportional policy & Unspecified & Architecture can become cap-driven \\
\bottomrule
\end{tabularx}
\end{table}

\subsection{Architecture definitions}
\begin{table}[htbp]
\centering
\caption{Dataset, architecture, and curriculum definitions used in the replication.}
\label{tab:architecture-definitions}
\small
\begin{tabularx}{\textwidth}{p{1.5cm}p{1.4cm}p{3.0cm}p{3.0cm}p{2.4cm}Y}
\toprule
Dataset & Input & Initial encoder & Mirrored decoder & Base classes & Incremental curriculum \\
\midrule
MNIST & $28\times28$ & $[200,100,75,20]$ & $[75,100,200,784]$ & $[1,7]$ & $[0,2,3,4,$ $5,6,8,9]$ \\
SD-19 & $28\times28$ & $[1000,500,250,50]$ & $[250,500,1000,784]$ & digits $0$--$9$ & uppercase and lowercase letters; feasibility subset in exported results \\
\bottomrule
\end{tabularx}
\end{table}

\subsection{Capacity matching}
\begin{draftpoints}
  \item conventional learner is initialized or retrained at the exact final encoder width reached by its paired NDL run, preventing a larger endpoint network from being credited as a growth-mechanism advantage
  \item Capacity matching is performed within replay regime and seed wherever endpoint widths vary; the corrected no-replay comparison uses each NDL seed's exact endpoint
  \item cumulative-cap reference is retained only as an architecture-shape comparison because its width is imposed by growth allowances rather than demonstrated to emerge organically
\end{draftpoints}
